\documentclass[11pt,a4paper]{article}
\usepackage[T1]{fontenc}
\usepackage{lmodern,amsmath,amssymb,amsthm,booktabs,array,microtype}
\usepackage[margin=25mm]{geometry}
\usepackage{xcolor}
\usepackage[colorlinks=true,linkcolor=teal!65!black,citecolor=teal!65!black,urlcolor=teal!65!black]{hyperref}
\hypersetup{pdftitle={Positive Measures, Cubic Congruences, and the Harmonic Radius},pdfauthor={Keio and Shaitetsu}}
\newtheorem{theorem}{Theorem}[section]
\newtheorem{lemma}[theorem]{Lemma}
\newtheorem{proposition}[theorem]{Proposition}
\newtheorem{corollary}[theorem]{Corollary}
\theoremstyle{remark}\newtheorem{remark}[theorem]{Remark}
\newcommand{\vp}{v_p}\newcommand{\Zp}{\mathbb Z_p}
\newcommand{\rad}{\operatorname{rad}}\newcommand{\diam}{\operatorname{diam}}\newcommand{\dist}{\operatorname{dist}}
\numberwithin{equation}{section}
\allowdisplaybreaks[1]\setlength{\emergencystretch}{2em}
\title{Positive Measures, Cubic Congruences,\\and the Harmonic Radius}
\author{Keio \and Shaitetsu}
\date{15 November 2026\\\small Research manuscript for independent mathematical and priority review}
\begin{document}\maketitle
\begin{abstract}
We give complete arguments for three families of previously stated conjectures.
First, both parity subsequences of the Baxter numbers are Stieltjes moment
sequences. Explicit probability densities on $[0,64]$ are obtained from the
Haar trace distribution on $\mathrm{SU}(3)$. An opposite-slice inequality
and an exactly normalized singular integral establish the positivity and
total mass needed for the even subsequence. Second, we prove Bala's cubic
supercongruences in OEIS A365029 for the complete mixed-binomial family,
all primes $p\ge5$, and every positive index and prime-power level.
Finite summation by parts matches valuation losses with complete and
half residue-block cancellation. Third, the absolute-radius estimate
of VMRC [161], reproduced with its full proof and credited to that source,
implies the harmonic-radius conjecture of Deng, Tang and Zhang and its
complete equality classification. A local boundary estimate and degree
integrality give a sharp positive gap of $1/15$ outside the exceptional
families. The graph application is an explicit
consequence of an existing corpus argument. No claim of worldwide
priority, nor of simultaneous certification of all six research-success
conditions, is made.
\end{abstract}

\section{Exact targets and the auxiliary problems}
The substantive analytic question in the first part is the following:
for $X=2+2\operatorname{Re}\operatorname{tr}U$, with $U$ Haar distributed
on $\mathrm{SU}(3)$ and density $f$, prove
\[
 f(t)>f(-t)\quad(0<t<8),\qquad
 \int_0^8\frac{f(t)-f(-t)}t\,dt=1.
\]
The nontrivial estimate is an inequality between opposite slices of the
deltoid trace density. Solving this auxiliary problem supplies a positive
measure for the unchanged Baxter even subsequence, including its zeroth
moment. The exact walk and character identities are proved before the
transfer. Thus the target is not replaced by a newly defined sequence.

For the arithmetic part the auxiliary question is to obtain one missing
$p$-adic order in a weighted inverse-square sum, uniformly at every depth.
Its weights change at two types of residue boundaries. A finite
summation-by-parts proof compensates every loss at those boundaries.
The resulting theorem is applied to Bala's original binomial family,
without changing its summation range or parameters. Classical weighted
cancellation methods of Beukers and Straub are credited explicitly.

The graph part is a different kind of contribution: a short comparison
transfers an existing corpus estimate to a named inequality.
The full structural proof is included so that an unreviewed corpus
claim is not used as an unexplained assumption. This transfer is not
presented as a new structural method.

The sources identify the Baxter parity questions as P1.7.5e/o on a
catalogue page dated 16 June 2026; Bala's entry is dated 27 August 2023;
the harmonic-radius conjecture appears in a 2012 preprint. These dates
do not support a claim that this manuscript resolves a problem known
to have remained open for twenty years. Section~\ref{sec:status}
distinguishes the mathematical claims from the remaining historical
questions.
% Standalone section fragment. Requires amsmath, amssymb, amsthm, hyperref.
% Draft proof for independent review; no historical-priority assertion.
\section{Both parity subsequences of the Baxter numbers are Stieltjes moments}
\label{sec:baxter-parity}

Let $B_0=1$ and let $B_n$ be the number of Baxter permutations of length
$n$. Thus $(B_n)_{n\ge0}=(1,1,2,6,22,92,422,\ldots)$. For $n\ge1$ the
standard enumeration formula \cite{ChungGrahamHoggattKleiman1978} is
\begin{equation}\label{eq:baxter-count}
 B_n=\frac{2}{n(n+1)^2}\sum_{k=1}^n
 \binom{n+1}{k-1}\binom{n+1}{k}\binom{n+1}{k+1}.
\end{equation}
The two parity moment conjectures, attributed to Stoyan Dimitrov, Bishal
Deb, and Alan Sokal, are recorded as Problems P1.7.5e and P1.7.5o in
Deb's positivity-problem catalogue: $(B_{2n})_{n\ge0}$ and
$(B_{2n+1})_{n\ge0}$ should each be Stieltjes moment sequences.\footnote{
\url{https://bishaldeb.com/positivity/themes/SMSCHTP/pp-hankel.html},
catalogue last updated 16 June 2026 and inspected 15 November 2026.
The catalogue gives no earlier original date or paper for these two assertions. Its displayed enumeration formula has
$n+2$ in the denominator, inconsistent with the stated Baxter-permutation
count already at $n=1$. We use the intrinsically defined Baxter counts,
with their standard enumeration \eqref{eq:baxter-count}; we do not
interpret that typographical inconsistency as a different sequence or
as a counterexample to the conjecture.}

\begin{theorem}\label{thm:baxter-parity}
Both $(B_{2n})_{n\ge0}$ and $(B_{2n+1})_{n\ge0}$ are Stieltjes moment
sequences represented by probability measures on $[0,64]$. Their
representing measures have strictly positive densities throughout
$(0,64)$. Consequently every minor of their Hankel matrices, with
strictly increasing nonnegative row and column indices, is strictly
positive. After division of the $n$th term by $64^n$, both sequences
are determinate Hausdorff moment sequences.
\end{theorem}

The main estimate below compares opposite slices of the $\mathrm{SU}(3)$
trace distribution. It converts a signed spectral folding into a
positive measure. The odd subsequence requires only the standard
self-adjoint walk representation; the even subsequence also requires
this slice comparison and an exact normalization at zero energy.

\subsection{The Baxter walk and a compact-group moment representation}

Let $\mathcal H=\ell^2(\mathbb N_0^2)$, and let $e$ be the unit vector at
$(0,0)$. Let $A$ be the symmetric adjacency operator whose allowed
increments are
\[
 (1,0),\ (-1,0),\ (0,1),\ (0,-1),\ (1,-1),\ (-1,1),
\]
with two loops at every vertex; a step leaving the quadrant is forbidden.
Then $A$ is a bounded self-adjoint operator.

We first verify the exact indexing
\begin{equation}\label{eq:baxter-walk}
 \langle e,A^ne\rangle=B_{n+1}\qquad(n\ge0).
\end{equation}
Consider three walkers at distinct integer positions $0,1,2$. At each
time every walker moves either zero or one unit, and their positions
must remain strictly ordered. Their consecutive gaps, minus one, give
a quadrant walk with exactly the six increments and two distinct loops
above. A return to gaps $(0,0)$ after $n$ steps means that the endpoints
are $k,k+1,k+2$ for some $0\le k\le n$. The nonintersecting-path
lemma therefore counts these returns by
\[
 \sum_{k=0}^n\det\left[\binom n{k+i-j}\right]_{i,j=0}^2.
\]
(The transposed convention gives the same determinant.) The elementary
binomial determinant evaluation
\[
 \det\left[\binom n{k+i-j}\right]_{i,j=0}^2
 =\frac{2\binom{n+2}{k}\binom{n+2}{k+1}\binom{n+2}{k+2}}
 {(n+1)(n+2)^2}
\]
follows by inserting factorials and expanding the $3\times3$
determinant; it is valid also at $k=0,n$ with the usual zero convention.
Summing gives \eqref{eq:baxter-walk} by \eqref{eq:baxter-count}.
This walk interpretation is also recorded in the literature on Baxter
permutations and rectangulations \cite{AsinowskiCardinalFelsnerFusy2025}.\footnote{The primary manuscript at
\url{https://igm.univ-mlv.fr/~fusy/Articles/rectangulations.pdf}, p.~29,
gives the same step set and the length-$n-1$ indexing for $B_n$.}

Let $U$ be Haar distributed on $\mathrm{SU}(3)$ and put
\[
 Z=\operatorname{tr}U=u+iv,\qquad X=2+Z+\overline Z.
\]
Irreducible $\mathrm{SU}(3)$ characters are indexed by pairs
$(a,b)\in\mathbb N_0^2$ and form an orthonormal basis of the class
functions. The fundamental tensor-product rule is
\[
 \chi_{1,0}\chi_{a,b}
 =\chi_{a+1,b}+\chi_{a-1,b+1}+\chi_{a,b-1},
\]
with terms having a negative index omitted; conjugation gives the reverse
three steps. Hence multiplication by $2+\chi_{1,0}+\chi_{0,1}$ has the
matrix $A$ in this basis, and the constant character corresponds to $e$.
It follows that
\begin{equation}\label{eq:baxter-su3}
 B_{n+1}=\mathbb E[X^n]\qquad(n\ge0).
\end{equation}
This is an equality for all $n$, not an inference from initial terms.

\subsection{The trace density and its slice comparison}

Define
\begin{align}
 D(u,v)&=27-18(u^2+v^2)+8(u^3-3uv^2)-(u^2+v^2)^2\notag\\
 &= (3-u)^3(u+1)-2(u^2+12u+9)v^2-v^4.
 \label{eq:deltoid-D}
\end{align}
The invariant trace distribution was studied in \cite{GuralnikWarnockZemach1985}.
We derive here the exact density and normalization needed below:
\begin{equation}\label{eq:trace-density}
 \rho(u,v)=\frac1{2\pi^2}\sqrt{D(u,v)_+},\qquad t_+=\max(t,0).
\end{equation}
For completeness, write the eigenvalues as
$z_1=e^{i\theta_1}$, $z_2=e^{i\theta_2}$,
$z_3=e^{-i(\theta_1+\theta_2)}$. The Weyl density on this torus is
$|\Delta(z)|^2/(6(2\pi)^2)$. The trace map is generically six-to-one;
its absolute Jacobian is $|\Delta(z)|/2$, and
$|\Delta(z)|^2=D(u,v)$. Change of variables gives
\eqref{eq:trace-density}. The region $D\ge0$ is the usual filled deltoid;
its real coordinate ranges from $-3/2$ to $3$. In particular
$-1\le X\le8$.

Set
\[
 I(u)=\int_{\mathbb R}\sqrt{D(u,v)_+}\,dv,
 \qquad
 f(\lambda)=\frac1{4\pi^2}I\left(-1+\frac\lambda2\right).
\]
Then $f$ is the probability density of $X$, supported on $[-1,8]$.
The function $\sqrt{D_+}$ has compact support and is uniformly
H\"older continuous of exponent $1/2$ there, including its zero
extension. In particular $f$ is H\"older continuous of exponent $1/2$.

\begin{lemma}[Opposite-slice domination]\label{lem:deltoid-domination}
For $0\le s\le1/2$,
\begin{equation}\label{eq:deltoid-chain}
 I(-1+s)\ \ge\ (1-s)^{3/2}I(-1)
 \ \ge\ (1-2s)^3I(-1)\ \ge\ I(-1-s).
\end{equation}
For $s>0$, the resulting comparison $I(-1+s)>I(-1-s)$ is strict.
\end{lemma}
\begin{proof}
At $u=-1$, $D(-1,v)=v^2(4-v^2)$ and $I(-1)=16/3$.
For the positive slice, direct subtraction gives
\begin{align*}
 &D(-1+s,\sqrt{1-s}\,w)-(1-s)^2D(-1,w)\\
 &\hspace{12mm}=s\bigl((4-s)^3-(1-s)(16+2s)w^2\bigr).
\end{align*}
When $|w|\le2$ and $0\le s\le1/2$, the right side is at least
$s^2(8+20s-s^2)\ge0$. Restricting the integral to
$|v|\le2\sqrt{1-s}$ and changing variables proves the first inequality;
its comparison is strict when $s>0$.

For the negative slice write
\[
 c=2+10s-s^2,\qquad d=2(1-2s)^{3/2}.
\]
The two roots in $v^2$ of $D(-1-s,v)$ are $c-d,c+d$.
For $0\le s<1/2$, substitution $v^2=c+dt$ therefore gives
\[
 I(-1-s)=d^2\int_{-1}^1\frac{\sqrt{1-t^2}}{\sqrt{c+dt}}\,dt.
\]
Since $d\le2$, the minimum over $-1\le t\le1$ of
$c+dt-(2+2t)$ occurs at $t=1$. Convexity gives
$(1-2s)^{3/2}\ge1-3s$, and consequently
\[
 c+d-4\ge4s-s^2\ge0.
\]
Thus $c+dt\ge2+2t$. Comparing with the same integral at $s=0$
proves $I(-1-s)\le(d^2/4)I(-1)=(1-2s)^3I(-1)$.
At $s=1/2$ the negative slice has zero area, so the inequality persists.
Finally $1-s\ge(1-2s)^2$ on this interval; taking the power $3/2$
proves the middle inequality of \eqref{eq:deltoid-chain}.
\end{proof}

It follows in particular that
\begin{equation}\label{eq:fold-positive}
 f(r)-f(-r)>0\qquad(0<r<8),
\end{equation}
where $f(-r)=0$ for $r>1$.

\subsection{Exact mass of the folded measure}

\begin{lemma}\label{lem:baxter-inverse-mass}
The symmetrically interpreted integral satisfies
\begin{equation}\label{eq:inverse-mass}
 \int_0^8\frac{f(r)-f(-r)}r\,dr=1.
\end{equation}
\end{lemma}
\begin{proof}
For $0\le r<1$, the complete-symmetric-function character expansion is
\[
 \det(I+rU)^{-1}=\sum_{m\ge0}(-r)^m\chi_{m,0}(U).
\]
It converges uniformly and absolutely, because
$|\chi_{m,0}(U)|\le(m+1)(m+2)/2$. Character orthogonality gives
\begin{equation}\label{eq:haar-inverse}
 \mathbb E[\det(I+rU)^{-1}]=1.
\end{equation}
To justify the boundary limit explicitly, write, for $1/2\le r<1$,
\[
 \det(I+rU)=r(1+r)\bigl(u-u_r+i\varepsilon_r v\bigr),
 \quad u_r=-(r+r^{-1}-1),\quad
 \varepsilon_r=\frac{1-r}{1+r}.
\]
Here $u_r\to-1$. Choose a fixed $R$ containing the support after
translation by every such $u_r$. For each $v$, the integral of the real
part of the resolvent against $\rho(u,v)$ equals
\[
 \frac1{r(1+r)}\int_0^R
 \frac{h}{h^2+\varepsilon_r^2v^2}
 \bigl[\rho(u_r+h,v)-\rho(u_r-h,v)\bigr]\,dh.
\]
Uniform H\"older continuity bounds the absolute integrand by
$C h^{-1/2}$, independently of $v$ and $r$; the relevant $v$ interval
is fixed and bounded. Dominated convergence therefore yields the
limit
\[
 \frac12\int_{\mathbb R}\int_0^R
 \frac{\rho(-1+h,v)-\rho(-1-h,v)}h\,dh\,dv
 =\int_0^8\frac{f(t)-f(-t)}t\,dt.
\]
Taking real parts and limits in \eqref{eq:haar-inverse} proves
\eqref{eq:inverse-mass}. The same H\"older estimate ensures ordinary
absolute convergence of the paired integral at zero.
\end{proof}

\subsection{Positive measures for the two subsequences}

The probability measure of $X^2$ has density
\begin{equation}\label{eq:baxter-odd-density}
 g_{\rm odd}(t)=\frac{f(\sqrt t)+f(-\sqrt t)}{2\sqrt t},
 \qquad 0<t<64.
\end{equation}
By \eqref{eq:baxter-su3}, its $n$th moment is $B_{2n+1}$ for every
$n\ge0$.

For the even subsequence define
\begin{equation}\label{eq:baxter-even-density}
 g_{\rm even}(t)=\frac{f(\sqrt t)-f(-\sqrt t)}{2t},
 \qquad 0<t<64.
\end{equation}
Equation \eqref{eq:fold-positive} makes this density strictly positive;
its integral is one by Lemma~\ref{lem:baxter-inverse-mass}. For every
$n\ge1$, changing variables $t=\lambda^2$ gives
\begin{align*}
 \int_0^{64}t^ng_{\rm even}(t)\,dt
 &=\int_0^8\lambda^{2n-1}\bigl(f(\lambda)-f(-\lambda)\bigr)\,d\lambda\\
 &=\int_{-1}^8\lambda^{2n-1}f(\lambda)\,d\lambda
 =B_{2n}.
\end{align*}
Its zeroth moment is $1=B_0$, as required. This proves both full moment
assertions, including the initial index which cannot be dropped.

Finally, for either positive density $g$ and strictly increasing
nonnegative integer indices $i_1<\cdots<i_m$,
$j_1<\cdots<j_m$, Andr\'eief's identity gives
\[
 \det\left[\int t^{i_a+j_b}g(t)\,dt\right]_{a,b=1}^m
 =\frac1{m!}\int
 \det[t_b^{i_a}]_{a,b=1}^m
 \det[t_b^{j_a}]_{a,b=1}^m\prod_{b=1}^m g(t_b)\,dt_b>0.
\]
The generalized Vandermonde determinants have the same strict sign on
$0<t_1<\cdots<t_m<64$, and this region has positive measure. Compact
support gives moment determinacy; rescaling by $64$ proves the final
Hausdorff assertion of Theorem~\ref{thm:baxter-parity}.

\paragraph{Priority and scope.}
The walk interpretation, the character tensor-product rule, the Haar
trace density, and the spectral moment construction are classical
inputs. The substantive step supplied here is the explicit opposite-slice
comparison together with the zero-energy normalization needed for the
even subsequence. Targeted literature checks are recorded separately;
this draft does not infer historical novelty from the absence of a
search result. The conjecture source is the named catalogue entry on a
page last updated 16 June 2026; an earlier original proposal date has
not been established.

\section{The complete cubic tower for Bala's mixed-binomial family}


\subsection{The original family and the result}
For integers $A\ge2$ and $B\ge1$, put $a_{A,B}(0)=1$ and
\begin{equation}\label{A:family}
a_{A,B}(N)=\sum_{k=0}^{N} F_{A,B}(N,k),\qquad
F_{A,B}(N,k)=\binom{N+k-1}{k}^{\!A}\binom{2k-1}{N}^{\!B}
\quad(N\ge1).
\end{equation}
At $k=0$ we use $\binom{-1}{N}=(-1)^N$, so
$F_{A,B}(N,0)=(-1)^{NB}$. At positive $k$ all binomial upper indices
are nonnegative, with the usual convention that a binomial coefficient
is zero when its lower index exceeds its upper index.

The sequence A365029 is $a_{2,1}$. Bala's entry, dated 27 August 2023,
conjectures both $a_{2,1}(p-1)\equiv a_{2,1}(0)\pmod{p^3}$ and
\[
a_{2,1}(np^r)\equiv a_{2,1}(np^{r-1})\pmod{p^{3r}}
\quad(p\ge5\text{ prime};\ n,r\ge1).
\]
The same entry proposes both assertions for the unchanged family
\eqref{A:family}, with all positive integers $A,B$ such that $A\ge2$.

\begin{theorem}\label{A:main}
For every prime $p\ge5$, integers $A\ge2$, $B\ge1$, and positive
integer $N$ divisible by $p$, one has
\begin{equation}\label{A:strong}
a_{A,B}(N)\equiv a_{A,B}(N/p)\pmod{p^{3\vp(N)}}.
\end{equation}
Moreover,
\begin{equation}\label{A:boundary}
a_{A,B}(p-1)\equiv1\pmod{p^{A+B}}.
\end{equation}
Consequently both of Bala's proposed supercongruences hold for every
member of his stated family and all the original parameters.
\end{theorem}

All rational congruences below are in $\Zp$. We set $\vp(0)=+\infty$.

\subsection{Inverse-square suffix sums}
\begin{lemma}\label{A:suffix}
Let $p\ge5$, let $p\mid N$, and put $\alpha=\vp(N)$. Define
\[
H_N(k)=\sum_{\substack{k<j<N\\p\nmid j}}j^{-2}
\qquad(0\le k<N).
\]
Then $H_N(0)\in p^\alpha\Zp$. If $0\le s\le\alpha$ and either
$p^s\mid k$ or $p^s\mid(2k+1)$, then
\[
H_N(k)\in p^s\Zp.
\]
\end{lemma}
\begin{proof}
The assertion for $s=0$ is just integrality. For $s\ge1$, set $q=p^s$.
Multiplication by 2 permutes the units modulo $q$, so
\[
\sum_{\substack{1\le j<q\\p\nmid j}}j^{-2}
\equiv\frac14\sum_{\substack{1\le j<q\\p\nmid j}}j^{-2}\pmod q.
\]
Since $3/4$ is a unit, this complete sum vanishes modulo $q$.
The reflection $j\mapsto q-j$ identifies the lower and upper half
sums modulo $q$; hence each half sum also vanishes, since 2 is a unit.

The endpoint $N$ is a multiple of $q$. If $q\mid k$, the suffix
$k<j<N$ consists of complete residue blocks, with multiples of $p$
omitted. If $q\mid2k+1$, then $k\equiv(q-1)/2\pmod q$ and the
suffix consists of one upper half-block followed by complete blocks.
Every such block has inverse-square sum zero modulo $q$. Taking
$k=0$, $s=\alpha$ gives the first assertion.
\end{proof}

\subsection{The primitive indices: summation by parts}
\begin{lemma}\label{A:primitive}
Under the hypotheses of Theorem~\ref{A:main},
\[
\sum_{\substack{1\le k<N\\p\nmid k}}F_{A,B}(N,k)
\in p^{3\alpha}\Zp,\qquad \alpha=\vp(N).
\]
\end{lemma}
\begin{proof}
Write
\[
X_k=\binom{N+k-1}{k-1},\qquad Y_k=\binom{2k-1}{N}
\quad(1\le k<N).
\]
Both are integers. For $p\nmid k$,
\begin{equation}\label{A:primitive-factor}
F_{A,B}(N,k)=\frac{N^A}{k^A}X_k^A Y_k^B.
\end{equation}
If $A\ge3$, each term already has valuation at least $A\alpha\ge3\alpha$.
It remains to treat $A=2$.

Put $Q_k=X_k^2Y_k^B$ and define
$e(m)=\max\{0,\alpha-\vp(m)\}$ for positive integers $m$.
The identity
\[
X_{k+1}-X_k=\frac Nk X_k
\]
implies
\begin{equation}\label{A:Xincrement}
X_{k+1}^2-X_k^2\in p^{e(k)}\mathbb Z
\quad(1\le k\le N-2).
\end{equation}
When $\vp(k)<\alpha$, this follows from the displayed identity
and the integrality of $X_k$ and $X_{k+1}+X_k$. Otherwise the desired
statement is simply integrality.

For $m\ge1$, put $D_m=\binom{m-1}{N-1}$. Pascal's identity gives
\[
Y_{k+1}-Y_k=D_{2k}+D_{2k+1},\qquad
D_m=\frac Nm\binom mN\in p^{e(m)}\mathbb Z.
\]
These identities remain valid when $m<N$, since both sides then vanish.
Writing
\[
P_k=\sum_{i=0}^{B-1}Y_{k+1}^{B-1-i}Y_k^i\in\mathbb Z,
\]
we obtain the exact decomposition
\begin{equation}\label{A:Qincrement}
Q_{k+1}-Q_k=
 (X_{k+1}^2-X_k^2)Y_{k+1}^B
 +X_k^2P_kD_{2k}+X_k^2P_kD_{2k+1}.
\end{equation}
The three terms have valuations at least $e(k)$, $e(2k)$, and
$e(2k+1)$ respectively.

By Lemma~\ref{A:suffix}, multiplication by $H_N(k)$ raises each of
these three valuations to at least $\alpha$. For the first two terms
use $s=\min\{\alpha,\vp(k)\}=\min\{\alpha,\vp(2k)\}$;
for the last use $s=\min\{\alpha,\vp(2k+1)\}$.
Thus
\[
(Q_{k+1}-Q_k)H_N(k)\in p^\alpha\Zp.
\]
Finite summation by parts now gives
\begin{align*}
\sum_{\substack{1\le k<N\\p\nmid k}}\frac{Q_k}{k^2}
&=Q_1H_N(0)+\sum_{k=1}^{N-2}(Q_{k+1}-Q_k)H_N(k)\\
&\in p^\alpha\Zp.
\end{align*}
Multiplying by $N^2$ and using \eqref{A:primitive-factor} proves the lemma.
\end{proof}

\subsection{The indices divisible by the prime}
We use the classical Jacobsthal congruence in the valuation-sensitive
quotient form: for $p\ge5$ and integers $a>b>0$,
\begin{equation}\label{A:jac}
\frac{\binom{pa}{pb}}{\binom ab}
\in1+p^{3+\vp(ab(a-b))}\Zp.
\end{equation}
This is the positive-parameter case of \cite[Lemma 5.1]{A:Straub};
endpoint binomial quotients equal one.

\begin{lemma}\label{A:divisible}
For every $k$ with $0\le k\le N$ and $p\mid k$,
\[
F_{A,B}(N,k)\equiv F_{A,B}(N/p,k/p)\pmod{p^{3\alpha}}.
\]
\end{lemma}
\begin{proof}
At $k=0$ the two terms are equal because $p$ is odd. If $k>0$ and
$2k\le N$, both terms vanish. Suppose $2k>N$ and write $s=\vp(k)$.
The scale-invariant factorizations
\[
\binom{N+k-1}{k}=\frac{N}{N+k}\binom{N+k}{k},\qquad
\binom{2k-1}{N}=\frac{2k-N}{2k}\binom{2k}{N}
\]
show that the quotient of the two $F$ terms is $R_1^A R_2^B$, where
$R_1,R_2$ are the Jacobsthal quotients for the displayed unshifted
binomials. Their precisions are at least
\[
\alpha+s+\vp(N+k),\qquad
\alpha+s+\vp(2k-N),
\]
respectively. If $s<\alpha$, both equal $\alpha+2s$, so
$R_1^AR_2^B-1\in p^{\alpha+2s}\Zp$.
Furthermore,
\[
\binom{N/p+k/p-1}{k/p}
=\frac{N/p}{k/p}\binom{N/p+k/p-1}{k/p-1}
\]
has valuation at least $\alpha-s$. The lower-scale summand therefore
has valuation at least $A(\alpha-s)\ge2(\alpha-s)$.
Multiplication by the quotient error yields valuation at least
$2(\alpha-s)+\alpha+2s=3\alpha$.

If $s\ge\alpha$, both quotient precisions are at least $3\alpha$.
The lower-scale summand is an integer, so the same conclusion follows.
This argument includes $k=N$ and does not assume that the individual
rational prefactors in the factorizations are $p$-integral.
\end{proof}

\begin{proof}[Proof of Theorem~\ref{A:main}]
Split \eqref{A:family} into the terms with $p\mid k$ and those with
$p\nmid k$. Lemma~\ref{A:divisible} transfers the former sum to
$a_{A,B}(N/p)$, and Lemma~\ref{A:primitive} annihilates the latter modulo
$p^{3\alpha}$. This proves \eqref{A:strong}.
For $N=np^r$, one has $\alpha=r+\vp(n)\ge r$, giving every original
index and level, including $n$ divisible by $p$.

For \eqref{A:boundary}, the $k=0$ term is 1. A nonzero term with
$k\ge1$ must have $k\ge(p+1)/2$. For these $k\le p-1$,
\[
\vp\binom{p+k-2}{k}=1,\qquad
\vp\binom{2k-1}{p-1}=1.
\]
Indeed each numerator factorial has exactly one factor $p$, while
both denominator factorials have none. Thus every remaining nonzero
summand has valuation $A+B$. This proves \eqref{A:boundary}; since
$A+B\ge3$, the originally proposed boundary modulus follows as well.
\end{proof}

\subsection{Scope and antecedents}
Jacobsthal's theorem and summation by parts are classical. In particular,
Straub's weighted cancellation lemma \cite[Lemma 5.6]{A:Straub}, generalizing
Beukers \cite[Lemma 3]{A:Beukers}, combines complete residue-block cancellation
with a compatible digit recurrence; reciprocal-square weights also occur
in Straub's proof of his Theorem 3.2. Thus weighted reciprocal-square
cancellation is not itself claimed as a new general idea. The particular
arithmetic mechanism here is the matching of the three increment loci
$k$, $2k$, and $2k+1$ with complete or half residue suffixes. It produces
all prime-power levels without induction on a finite list of checked
levels. This is an additional valuation argument, not an identity alone.
The public proof portfolio \cite{A:Bajaj} had already recorded the
stronger boundary congruence and the first two adjacent levels for
A365029, and explicitly isolated the higher tower as unresolved.
Those earlier results are not claimed as new here. The present proof
is self-contained apart from the stated classical Jacobsthal theorem.
Its historical originality remains subject to specialist literature review.


\section{The harmonic index and graph radius}

\subsection{Statement, provenance, and the transfer}
Every graph below is finite, undirected, simple, and connected. Write
$n=|V(G)|$, $d_v=\deg_G(v)$, $\Delta=\max_v d_v$, and
\[
 H(G)=\sum_{uv\in E(G)}\frac{2}{d_u+d_v},\qquad
 M(G)=\sum_{uv\in E(G)}\frac1{\max(d_u,d_v)}.
\]
Empty sums for $K_1$ are zero. The vertex radius is
$r=\rad(G)=\min_{v\in V(G)}\max_{w\in V(G)}\dist_G(v,w)$.
Give every edge length one and allow centers anywhere in the resulting
metric graph $|G|$, with the original vertices as demand points:
\[
 \rho(G)=\min_{x\in|G|}\max_{w\in V(G)}\dist_{|G|}(x,w).
\]
Deng, Tang and Zhang posed the following unchanged conjecture in
\cite[Conjecture~6.1]{H:DTZ}: for every connected graph except an even path,
$H(G)\ge\rad(G)$. Shi \cite[Conjecture~1.4]{H:Shi} restated this problem
and proved it for sufficiently dense triangle-free graphs; the original
paper treated trees and unicyclic graphs and gave a cyclomatic-number bound.

\begin{theorem}\label{H:thm:harmonic}
For a finite connected simple graph with $\Delta\ge3$,
\begin{equation}\label{H:eq:harmonicbound}
 H(G)-\rad(G)\ge\frac{\Delta-3}{2(\Delta+1)}.
\end{equation}
For all finite connected simple graphs, $H(G)<\rad(G)$ holds precisely
for the paths $P_{2m}$ with $m\ge2$. Equality holds precisely for
$K_1$, $K_2$, and the even cycles. Every other graph has $H(G)>\rad(G)$.
\end{theorem}

\begin{theorem}[Absolute-radius inequality, VMRC {\rm[161]}]\label{H:thm:abs}
Every finite connected simple graph satisfies $M(G)\ge\rho(G)$.
\end{theorem}
For self-containment, Sections~\ref{H:sec:metric}--\ref{H:sec:sum} reproduce
and check the full proof from \cite{H:VMRC161}; no part of that theorem
is assumed as an opaque import. The essential new estimate in that
source is the degree-threshold inequality of Proposition~\ref{H:prop:cut}.
Its reductions develop the earlier bag method of Cygan, Pilipczuk and
\v{S}krekovski \cite{H:CPS}. We credit that estimate and all its reductions
to their existing sources. The new conclusion of this note comes from
the following comparison, rather than from a new proof of the
absolute-radius theorem.

\begin{proposition}[Harmonic defect transfer]\label{H:prop:htransfer}
For every connected graph of order at least three,
\begin{equation}\label{H:eq:sharptransfer}
 H(G)-r\ge\frac{n-2r-(\Delta-1)(r-\rho(G))}{\Delta+1}.
\end{equation}
\end{proposition}
\begin{proof}
For each edge write $d\le e$ for its endpoint degrees, and put
$q_{uv}=2/(d+e)-1/e=(e-d)/[e(d+e)]\ge0$. The contribution to the
degree-counting identity is
\[
 \frac1d+\frac1e-\frac4{d+e}
 =\frac{(e-d)^2}{de(d+e)}
 =\left(\frac ed-1\right)q_{uv}.
\]
Summing $1/d_u+1/d_v$ over the edges counts every vertex exactly once.
Consequently, recovering the identity of Shi \cite[Theorem~3.3]{H:Shi},
\[
 n-2H(G)
 =\sum_{uv}\left(\frac ed-1\right)q_{uv}
 \le(\Delta-1)\sum_{uv}q_{uv}
 =(\Delta-1)(H(G)-M(G)).
\]
Thus $(\Delta+1)H(G)\ge n+(\Delta-1)M(G)
\ge n+(\Delta-1)\rho(G)$ by Theorem~\ref{H:thm:abs}.
Subtracting $(\Delta+1)r$ proves the assertion.
\end{proof}

\begin{proof}[Proof of Theorem~\ref{H:thm:harmonic}]
Lemmas~\ref{H:lem:half} and \ref{H:lem:tree}, proved below, give
$0\le r-\rho(G)\le1/2$ and $n-2r\ge\Delta-2$.
Substitution in \eqref{H:eq:sharptransfer} gives
\[
 H(G)-r\ge\frac{\Delta-2-(\Delta-1)/2}{\Delta+1}
 =\frac{\Delta-3}{2(\Delta+1)}.
\]
This is strictly positive when $\Delta\ge4$.
Suppose $\Delta=3$ and $H(G)=r$. Equality must hold in every
inequality used above. In particular $n-2r=1$, so $n$ is odd.
For every edge whose degrees differ, equality in the defect comparison
requires $e/d-1=2$, forcing $(d,e)=(1,3)$. There can be no vertex of
degree two: a shortest path from such a vertex to a degree-three vertex
would contain an edge with endpoint degrees $(2,1)$ or $(2,3)$.
Hence all degrees are one or three. The handshaking identity then forces
$n$ to be even, a contradiction. Thus $H(G)>r$ also when $\Delta=3$.

A connected graph with maximum degree two is a path or a cycle.
For $n\ge3$,
\[
 H(P_n)=\frac43+\frac{n-3}{2}=\frac n2-\frac16,
 \qquad H(C_n)=\frac n2,
 \qquad \rad(P_n)=\rad(C_n)=\lfloor n/2\rfloor.
\]
These formulas give all stated strict inequalities and cycle equalities.
Finally $H(K_1)=\rad(K_1)=0$ and $H(K_2)=\rad(K_2)=1$.
\end{proof}

\subsection{Metric centers and spanning trees}\label{H:sec:metric}
Distances between vertices in $|G|$ agree with graph distances. With
finitely many vertices, the maximum of their distance functions is a
continuous function on the compact space $|G|$, so an absolute center
exists.
\begin{lemma}\label{H:lem:half}
An absolute center can be chosen at a vertex or an edge midpoint.
In particular $\rad(G)-1/2\le\rho(G)\le\rad(G)$.
\end{lemma}
\begin{proof}
Consider an edge $uv$ with coordinate $t\in[0,1]$ measured from $u$.
For each vertex $w$ the distance from $t$ to $w$ is
\[
 \min\{\dist_G(u,w)+t,\ \dist_G(v,w)+1-t\}.
\]
Both affine functions have integer intercepts and slopes $\pm1$. Two
such functions of opposite slopes can meet only at a half-integer.
Consequently every individual distance, and their maximum, is affine
on each of the open half-edges. Its minimum on the closed edge is
attained at an endpoint or midpoint. Distances from a center so chosen
to vertices are multiples of $1/2$. Rounding any absolute center to a
nearest vertex increases all its distances by at most $1/2$. The vertex
centers are allowed when minimizing for $\rho$.
\end{proof}
\begin{lemma}\label{H:lem:tree}
For a tree $T$ of diameter $L$, one has $\rho(T)=L/2$ and
$\rad(T)=\lceil L/2\rceil$. If a connected graph $H$ on $m$ vertices
has a vertex with at least $j\ge2$ incident edges, then it admits a
spanning tree of diameter at most $m-j+1$. In particular
$2\rad(H)\le m-j+2$.
\end{lemma}
\begin{proof}
Let a longest path in $T$ have endpoints $a,b$. For any vertex $w$, its
branch meets the path at a unique point at distance $s$ from $a$; write
$t$ for the distance from that point to $w$. The two paths to the
endpoints give $t\le\min(s,L-s)$. Thus the distance from the midpoint
to $w$ is at most $|s-L/2|+\min(s,L-s)=L/2$. The endpoints force a
radius of at least $L/2$ by the triangle inequality. Rounding the
midpoint yields the asserted integer radius.

The $j$ chosen incident edges form a tree. Extend this tree, one edge
to a new vertex at a time, to a spanning tree $T$ of $H$. At least $j$
components of $T-v$ meet the chosen vertex $v$ by distinct edges.
Each contains a leaf of $T$, as a farthest vertex of that component
from $v$ is a leaf. A longest path in $T$ uses at most two leaves. It
therefore uses at most $m-(j-2)$ vertices and has at most $m-j+1$ edges.
Distances in $H$ can only be shorter than those in its spanning tree.
The vertex-radius bound now follows from $2\lceil L/2\rceil\le L+1$.
\end{proof}

\subsection{Reductions of a minimal counterexample}\label{H:sec:red}
Suppose Theorem~\ref{H:thm:abs} fails. Among counterexamples $G$ choose
one minimizing the pair $(\beta(G),|E(G)|)$ lexicographically, where
$\beta(G)=|E(G)|-|V(G)|+1$ is the cyclomatic number. This order on
pairs of nonnegative integers is well founded. A graph smaller in this
order satisfies $M\ge\rho$. Stars with at least two leaves have
$M=\rho=1$. On a path of order $m\ge3$, each edge contributes $1/2$
to $M$ and Lemma~\ref{H:lem:tree} gives $M=\rho=(m-1)/2$. Also
$(M,\rho)=(0,0)$ for $K_1$ and $(1,1/2)$ for $K_2$.

\begin{lemma}\label{H:lem:delete}
If $e$ is a nonbridge edge of $G$, then $M(G-e)>M(G)$. If $z$ is
a leaf with neighbor $v$ of degree $d\ge2$, and $b$ of the remaining
neighbors of $v$ have degree less than $d$, then
\begin{equation}\label{H:eq:oneleaf}
 M(G)-M(G-z)=\frac1d-\frac{b}{d(d-1)}\ge0,\qquad
 0\le\rho(G)-\rho(G-z)\le\tfrac12.
\end{equation}
The leaf formulas apply to any connected graph in the stated range.
\end{lemma}
\begin{proof}
Deleting a nonbridge preserves the demand vertices, removes potential
center locations, and cannot shorten distances. It reduces $\beta$ by
one. Hence $M(G-e)\ge\rho(G-e)\ge\rho(G)>M(G)$.

Deleting $z$ removes $1/d$ from $M$. If the other endpoint of an edge
at $v$ has degree less than $d$, its weight increases from $1/d$ to
$1/(d-1)$. The other weights stay fixed. As $b\le d-1$, this proves
the index formula.

Project a center lying on the removed pendant interval to $v$. All
distances to retained vertices weakly decrease, proving the first
radius inequality. Conversely, choose a vertex or midpoint center $y$
of $G-z$, and put $r_0=\rho(G-z)$. If $\dist(y,v)\le r_0-1/2$,
then $\dist(y,z)\le r_0+1/2$ after restoration. Otherwise
$\dist(y,v)=r_0$ by Lemma~\ref{H:lem:half}. As $G-z$ has at least two
vertices, $r_0\ge1/2$. Move $y$ a distance $1/2$ along a shortest
path toward $v$. The distances to retained vertices increase by at
most $1/2$, and the distance to $z$ is $r_0+1/2$.
\end{proof}

A vertex adjacent to a leaf will be called its support. A support of
degree two in $G$ has a second neighbor of degree at least two, because
$G$ is not a star. Formula~\eqref{H:eq:oneleaf} then gives a drop in
$M$ of $1/2$, at least the drop in $\rho$. Deleting this leaf would
produce a smaller counterexample with the same $\beta$. Thus every
support in $G$ has degree at least three.

If $z,z'$ are leaves with the same support $v$, delete $z$. In the
remaining graph a center can be chosen outside the pendant interval
leading to $z'$, except possibly at $v$. Indeed, projecting a point
of that interval to $v$ decreases distances to all other vertices,
and puts $z'$ at distance one. Any point on that interval has
eccentricity at least one, since another neighbor of $v$ remains.
Restoring $z$ now gives $z$ the same distance from the center as $z'$.
The radius is unchanged and $M$ weakly decreases upon deletion.
Minimality gives a contradiction. Each support therefore has exactly
one leaf.

\begin{lemma}[Splitting at a cut vertex]\label{H:lem:split}
For every vertex $v$ of the minimal counterexample $G$, at most one
component of $G-v$ contains an edge.
\end{lemma}
\begin{proof}
Otherwise partition $V(G)\setminus\{v\}$ into unions $A,B$ of components
of $G-v$, each containing an edge. For $i\in\{A,B\}$ let $t_i$ be
the number of neighbors of $v$ in $i$, and $m_i$ the number of edges
within $i$. All four numbers are positive. Form $G_A$ by retaining
$G[A\cup\{v\}]$ and adding $t_B$ new leaves at $v$. Define $G_B$
symmetrically. These graphs are connected, $|E(G_A)|=|E(G)|-m_B$,
and $\beta(G_A)=m_A+t_A-|A|$. The analogous expression for $G_B$
gives $\beta(G_A)+\beta(G_B)=\beta(G)$. Both graphs are therefore
smaller. Each retained edge has its original endpoint degrees.
The $t_A+t_B=d_v$ artificial edges across the two graphs each have
weight $1/d_v$, which gives
\begin{equation}\label{H:eq:splitindex}
 M(G_A)+M(G_B)=M(G)+1 .
\end{equation}

Set $r_i=\rho(G_i)$. In each $G_i$, an artificial leaf and a retained
neighbor of $v$ are at distance two, so $r_i\ge1$. Choose a center
$y_i$ in the retained part: if a center lies on an artificial pendant
edge, project it to $v$. This weakly decreases its distances to all
vertices other than the tip of that edge, and places that tip at
distance $1\le r_i$. If $s_i=\dist(y_i,v)$, the distance to an
artificial leaf is $s_i+1$, so $s_i\le r_i-1$.

In $|G|$ join $y_A$ to $v$ and $v$ to $y_B$ by shortest paths. Their
concatenation is a geodesic since the two retained parts meet only at
$v$. With signed distance $u$ from $v$ along this path, its ends have
coordinates $-s_A,s_B$. Via $y_A$ and $y_B$, respectively, the
distances to the original vertices on the two sides (including $v$)
are at most
\[
 L(u)=r_A+s_A+u,\qquad H(u)=r_B+s_B-u .
\]
If these functions meet on the interval, their common value is
$(r_A+s_A+r_B+s_B)/2\le r_A+r_B-1$. If $L$ dominates throughout,
at $u=-s_A$ their maximum is $r_A$. If $H$ dominates throughout,
at $u=s_B$ the maximum is $r_B$. Both are at most $r_A+r_B-1$.
Consequently $\rho(G)\le r_A+r_B-1\le M(G_A)+M(G_B)-1=M(G)$ by
minimality and \eqref{H:eq:splitindex}, a contradiction.
\end{proof}

\begin{proposition}\label{H:prop:core}
The minimal counterexample consists of a two-connected induced subgraph
$K$ of order at least three together with at most one leaf attached to
each vertex of $K$. Every edge of $K$ is a nonbridge in $G$.
\end{proposition}
\begin{proof}
A tree containing a path of length four violates Lemma~\ref{H:lem:split}
at the middle of that path. A tree of diameter at most two is a
singleton, an edge or a star. In a tree of diameter three, every vertex
lies on the central edge of a diametral path or is adjacent to one of
its ends: the midpoint covers the tree within $3/2$ by
Lemma~\ref{H:lem:tree}. Each central vertex supports at most one leaf,
so this tree is a path. These possibilities were excluded above.

Thus $G$ contains a cycle. Its vertices induce a two-connected graph,
where two-connected means having at least three vertices and remaining
connected after deletion of any one vertex. Extend that set to a
maximal vertex set inducing a two-connected subgraph $K$. If a component
outside $K$ has two distinct attachment vertices in $K$, it gives a
path between them with all interior vertices outside $K$. Adding the
path preserves two-connectivity: deletion of an interior vertex leaves
each remaining piece attached at an end, and deletion of a vertex of
$K$ leaves $K$ connected with at least one end of the path attached.
This contradicts maximality. Each outside component therefore has a
unique attachment $v$. If it contained an edge, deleting $v$ would
separate that edge from $K-v$, which is connected and has at least
two vertices. Lemma~\ref{H:lem:split} excludes this. Each outside
component is a single leaf, and the previous reduction bounds their
number per support. Finally, every vertex of $K$ has at least two
neighbors in $K$. For a core edge $uv$, choose another core neighbor
$w$ of $u$; a path in $K-u$ from $w$ to $v$ completes an alternative
route from $u$ to $v$. Hence $uv$ is a nonbridge.
\end{proof}

\subsection{A bound for each degree cut}\label{H:sec:cuts}
Fix the subgraph $K$ of Proposition~\ref{H:prop:core} and write
$n=|V(K)|$. Degrees $d_v$ in this section are always measured in $G$.
At a core vertex $v$, let $p_v$ be $1$ if $v$ supports a leaf and
$0$ otherwise. Define
\[
 \ell_v=\#\{w\sim v:d_w<d_v\},\quad h_v=d_v-\ell_v,\quad
 g_v=d_v(d_v-1),\quad a_v=\ell_v/g_v,\quad b_v=h_v/g_v .
\]
In particular $\ell_v$ includes the leaf when $p_v=1$.

The next lemma includes the local-minimum and minimum-degree properties
of bags from Lemma~4.2 of \cite{H:CPS}.
\begin{lemma}\label{H:lem:local}
Every core vertex has $d_v\ge3$ and $\ell_v\ge1$. For a core edge $uv$
with $d_u<d_v$ we have $a_u>b_v$.
\end{lemma}
\begin{proof}
On deleting $uv$, each remaining edge at an endpoint $x$ gains $1/g_x$
precisely when its other endpoint has degree strictly below $d_x$.
No remaining edge occurs at both $u$ and $v$, since $G$ is simple.
Subtracting the weight of the deleted edge gives
\begin{align}
 M(G-uv)-M(G)
 &=a_u+\frac{\ell_v-1}{g_v}-\frac1{d_v}
   =a_u-b_v &&(d_u<d_v),\label{H:eq:strict}\\
 &=\frac{\ell_u+\ell_v}{d(d-1)}-\frac1d
   &&(d_u=d_v=d).\label{H:eq:equal}
\end{align}
Both changes must be positive by Lemma~\ref{H:lem:delete}. If
$\ell_u=0$, there is no leaf at $u$. A core neighbor of higher degree
would contradict \eqref{H:eq:strict}, since $b_v\ge0$. A neighbor of
equal degree would have $\ell_v\le d-1$, contradicting
\eqref{H:eq:equal}. These exhaust the neighbors of $u$. Finally, a
vertex of full degree two already has two neighbors in $K$; all
vertices of $K$ have degree at least two. It could have no lower-degree
neighbor.
\end{proof}

A core vertex of minimum full degree must have a leaf, by
Lemma~\ref{H:lem:local}. Thus $K$ is smaller than $G$ in the order of
Section~\ref{H:sec:red}: removing the leaves leaves $\beta$ unchanged
and reduces the edge count. Let $\delta=M(G)-M(K)$. Repeated leaf
deletion and minimality give
\begin{equation}\label{H:eq:delta}
 0\le\delta<\rho(G)-\rho(K)\le1 .
\end{equation}
The last bound holds because a center of $K$ reaches each added leaf
within $\rho(K)+1$. If $\Delta=3$, every core vertex has full degree
three. Its only possible lower-degree neighbor is its leaf, giving
$\ell_v=1$ at every core vertex. Equation~\eqref{H:eq:equal} gives a
zero change at every core edge, a contradiction. Hence $\Delta\ge4$.

\begin{proposition}[Degree-cut estimate]\label{H:prop:cut}
For an integer $3\le k<\Delta$ let $S_k=\{v\in V(K):d_v\le k\}$,
and let $N_k$ count edges of $K$ with exactly one end in $S_k$. Then
\begin{equation}\label{H:eq:cut}
 N_k\le k(k-1)\delta .
\end{equation}
\end{proposition}
\begin{proof}
Write $S=S_k$. Orient each unequal-degree core edge from its lower
endpoint to its higher endpoint. Denote by $I$ the directed edges
inside $S$ and by $F$ the directed edges leaving $S$. Every crossing
edge leaves $S$. Let $e_v$ count the core neighbors with degree $d_v$.
At $v$ there are $\ell_v-p_v$ incoming and $h_v-e_v$ outgoing edges.
Since $h_v a_v=\ell_v b_v$, we have
\[
 p_v b_v=(h_v-e_v)a_v-(\ell_v-p_v)b_v+e_v a_v .
\]
Sum this identity over $v\in S$. An internal directed edge $u\to v$
contributes $a_u-b_v$, and an edge leaving $S$ contributes $a_u$:
\begin{equation}\label{H:eq:flow}
 \sum_{v\in S}p_v b_v
 =\sum_{(u,v)\in I}(a_u-b_v)+\sum_{(u,v)\in F}a_u
        +\sum_{v\in S}e_v a_v .
\end{equation}
Each term of the first sum is positive by Lemma~\ref{H:lem:local}.

Delete precisely the leaves whose support belongs to $S$ and call the
result $G_S$. To compute the change of $M$, first consider a deleted
leaf at $v$. It has weight $1/d_v$, and the $\ell_v-1$ edges to
strictly lower core neighbors each gain $1/g_v$. Indeed, on a strict
edge with original degrees $d<e$, the new maximum remains $e$ when
the higher endpoint keeps its leaf or has none. If the higher endpoint
loses a leaf, the new maximum is $e-1$, as $d\le e-1$, irrespective
of any deletion at the lower endpoint. Thus the loss from that leaf,
after these gains, is $1/d_v-(\ell_v-1)/g_v=b_v$.

On a core edge of equal degrees $d$, the maximum decreases precisely
when both endpoints lose a leaf. In that case the edge gains
$1/[d(d-1)]$. Let $Q$ consist of these unordered core edges. The
exact total change is therefore
\begin{equation}\label{H:eq:subset}
 M(G)-M(G_S)=\sum_{v\in S}p_v b_v-\sum_{uv\in Q}\frac1{g_u}.
\end{equation}
For $uv\in Q$, the last sum of \eqref{H:eq:flow} counts $a_u+a_v$
on that edge. This is at least $2/g_u$ since $g_u=g_v$ and
$\ell_u,\ell_v\ge1$. All other contributions to that last sum are
nonnegative. Subtracting the correction in \eqref{H:eq:subset} from
\eqref{H:eq:flow}, and then deleting the remaining leaves one at a time
by Lemma~\ref{H:lem:delete}, gives
\[
 \sum_{(u,v)\in F}a_u\le M(G)-M(G_S)\le M(G)-M(K)=\delta .
\]
For every $(u,v)\in F$, one has $\ell_u\ge1$ and $d_u\le k$, so that
$a_u\ge1/[k(k-1)]$. There are $N_k$ such edges, giving \eqref{H:eq:cut}.
\end{proof}

\subsection{Summing the cuts and bounding the absolute radius}\label{H:sec:sum}
\begin{proof}[Completion of the proof of Theorem~\ref{H:thm:abs}]
Continue with the hypothetical minimal counterexample. Write
$t=\sum_{v\in V(K)}p_v$, so $G$ has $n+t$ vertices. In a graph
without isolated vertices, the sum of $1/d_u+1/d_v$ over the edges
counts each vertex once. Subtracting $2M(G)$ gives
\[
 n+t-2M(G)=\sum_{uv\in E(G)}\left|\frac1{d_u}-\frac1{d_v}\right|.
\]
A leaf at $v$ contributes $1-1/d_v$. A core edge with degrees
$d<e$ contributes $\sum_{k=d}^{e-1}1/[k(k+1)]$, counting exactly
the cuts it crosses. Equal-degree edges contribute zero. We obtain
\begin{equation}\label{H:eq:sum}
 n-2M(G)=\sum_{k=3}^{\Delta-1}\frac{N_k}{k(k+1)}
                     -\sum_{v\in V(K)}\frac{p_v}{d_v}.
\end{equation}
When all leaves are deleted, the retained core-edge weights can only
increase. Thus the sum of the removed weights $\sum_v p_v/d_v$ is at
least $\delta$. By Proposition~\ref{H:prop:cut},
\begin{align}
 n-2M(G)
 &\le \left(\sum_{k=3}^{\Delta-1}\frac{k-1}{k+1}-1\right)\delta
 \le(\Delta-4)\delta
 \le\Delta-4 . \label{H:eq:lowerM}
\end{align}
There are $\Delta-3$ summands, each at most one; the last two
inequalities use $\delta\ge0$, $\Delta\ge4$ and \eqref{H:eq:delta}.

Choose a core vertex of full degree $\Delta$. It has at least
$\Delta-1$ neighbors in $K$, since at most one is a leaf. By
Lemma~\ref{H:lem:tree}, $K$ has a spanning tree of diameter at most
$n-(\Delta-1)+1=n-\Delta+2$. Restore each original leaf and its
incident edge to that tree. Between any two vertices of the resulting
tree $T$, the path consists of a core-tree path and at most two leaf
edges. Thus $\diam(T)\le n-\Delta+4$. Distances in $G$ from a
diametral midpoint of $T$ are at most their distances in $T$. Applying
Lemma~\ref{H:lem:tree} and \eqref{H:eq:lowerM} gives
\begin{equation}\label{H:eq:contradiction}
 2\rho(G)\le\diam(T)\le n-\Delta+4\le 2M(G).
\end{equation}
This contradicts $M(G)<\rho(G)$ and proves Theorem~\ref{H:thm:abs}.
\end{proof}


\subsection{A sharp positive gap}
The preceding transfer can be sharpened by a local boundary argument
when the maximum degree is three. The conclusion below is additional
to the coarse transfer; no assertion of historical priority is made.

\begin{theorem}\label{H:thm:sharp-gap}
If $G$ is a finite connected simple graph other than a path of even
order, an even cycle, or $K_1$, then
\[
 H(G)-\rad(G)\ge\frac1{15}.
\]
The constant $1/15$ is best possible and is attained at infinitely many
orders.
\end{theorem}
\begin{proof}
For maximum degree $\Delta\ge4$, \eqref{H:eq:harmonicbound} gives a gap
of at least $1/10$. Odd paths and odd cycles have gaps $1/3$ and $1/2$,
respectively. It remains to consider $\Delta=3$.

Put $r=\rad(G)$, $\eta=r-\rho(G)$ and
$q_{uv}=(e-d)/[e(d+e)]$ for edge degrees $d\le e$. Retaining the slack
in Proposition~\ref{H:prop:htransfer} gives the exact identity
\begin{equation}\label{H:eq:gap-slack}
 4(H(G)-r)=n-2r-2\eta+2(M(G)-\rho(G))+S,
 \qquad
 S=\sum_{uv\in E(G)}\left(3-\frac ed\right)q_{uv}\ge0.
\end{equation}
We know $n-2r\ge1$, $2\eta\le1$, and $M(G)\ge\rho(G)$.
If $n-2r\ge2$, the right side is at least one, so the gap is at least
$1/4$. Otherwise $n-2r=1$, and $n$ is odd. There must be a degree-two
vertex, since a graph all of whose degrees are odd has even order.

Consider a component of the subgraph induced by degree-two vertices.
It has maximum degree two and cannot be a cycle, since a cycle component
would have no edges leaving it and would be the whole connected graph.
It is therefore a path, allowing a single vertex as a path. Its vertices
have exactly two edges in total to vertices outside that component.
Those outside degrees are one or three. Both cannot be one: in that
case the path and its two leaves would form the whole graph, contrary
to $\Delta=3$. Thus the two boundary edges have types $(2,3),(2,3)$,
or types $(1,2),(2,3)$. Their respective contributions to $S$ are
\[
 \left(3-\frac32\right)\frac1{15}=\frac1{10},
 \qquad
 (3-2)\frac16=\frac16.
\]
Hence $S\ge1/5$. Equation~\eqref{H:eq:gap-slack} now gives
$H(G)-r\ge1/20$. Every possible harmonic edge weight for maximum
degree three belongs to $(1/30)\mathbb Z$, because the endpoint-degree
sums are among $2,3,4,5,6$. Since $r$ is integral, the positive gap is
a multiple of $1/30$, and being at least $1/20$ it is at least $1/15$.

For sharpness, let $T_m$ consist of a vertex $v$ with two leaf neighbors
and a third arm that is a path of length $2m$ from $v$ to a leaf,
where $m\ge1$. This tree has order $2m+3$ and diameter $2m+1$, so
$\rad(T_m)=m+1$. Its two short leaf edges contribute one in total;
the first and last edges of the long arm contribute $2/5$ and $2/3$;
its remaining $2m-2$ edges contribute $1/2$ each. Thus
\[
 H(T_m)=1+\frac25+\frac23+\frac{2m-2}{2}
 =m+\frac{16}{15},
 \qquad H(T_m)-\rad(T_m)=\frac1{15}.
\]
This proves both the bound and sharpness.
\end{proof}

\subsection{Scope of the conclusion and novelty limits}
The proof above has no numerical or bounded-order hypothesis, and proves
the original harmonic-radius statement with its original index and
radius. It also determines every equality and failure case.
The proof does not establish the separate Fajtlowicz average-distance
conjecture $\mu(G)\le R(G)$.

The corpus already contains the absolute-radius theorem and its
substantive degree-cut argument. The additional harmonic comparison is
elementary. Therefore the present note does not certify a wholly new
methodology, and it must not be described as satisfying a requirement
that rules out direct consequences of an existing corpus theorem.
Searches of the corpus, primary web sources, and SciSpace on
15 November 2026 located the original conjecture and partial results,
but no earlier unrestricted harmonic-radius theorem.
This is a literature-search outcome, not a proof of historical priority
or a guarantee that no earlier solution exists.


\section{Verification, antecedents, and the six success conditions}
\label{sec:status}
The proofs have been checked independently within this research session.
The Baxter argument was checked at its three essential interfaces:
the exact walk-to-character identification, the opposite-slice inequality,
and the boundary-limit normalization of the even measure. The arithmetic
argument was checked at every valuation stratum and endpoint, including
$p\mid n$. Its 135 exact supplementary tests used $p\in\{5,7,11\}$,
$1\le r,n\le3$, and five exponent pairs, with largest summation index
$3993$. The graph argument and its copied structural input were reviewed
in full; exact rational computations on all 996 connected graphs of
orders one through seven agree with the theorem. None of these finite
computations is used to replace a quantified step in a proof. No proof
assistant formalization or external specialist endorsement is asserted.

\begin{center}\small
\begin{tabular}{@{}p{0.08\textwidth}p{0.84\textwidth}@{}}
\toprule
Gate & Assessment in this manuscript\\\midrule
G1 & The named source and original statement are identified for each target.
The auxiliary slice and residue questions are not substituted for them.\\[3pt]
G2 & Original sequences, moduli, indices, and graph invariants are retained.
The Baxter catalogue's inconsistent displayed formula is documented;
the target is its explicitly defined Baxter-permutation count.\\[3pt]
G3 & The inspected sources present the assertions as conjectures, and the
closest located arithmetic attempt leaves higher levels open. This is
evidence from inspected sources, not an exhaustive historical certificate.\\[3pt]
G4 & Complete unconditional proofs of the stated mathematical assertions
are supplied, using the classical results explicitly identified in the text.\\[3pt]
G5 & The Baxter slice comparison and the arithmetic boundary estimate do
actual positivity and valuation work. Their historical originality remains
to be reviewed. The graph application is a consequence of VMRC [161]
and does not satisfy a demand for an independent new structural method.\\[3pt]
G6 & Not certified. Targeted primary-source and scholarly searches located
no identical complete resolution, but absence of a search hit does not
establish absence of prior literature.\\\bottomrule
\end{tabular}
\end{center}

The closest arithmetic predecessor inspected is the public proof portfolio
\cite{A:Bajaj}, at commit
\texttt{e6955c6a329a713e7ed7b3f1dca165f587131c90}.
It proves the boundary congruence and the first two adjacent levels for
$A=2,B=1$, and labels its final higher-level digit-cancellation statement
unproved. The boundary result is therefore reproduced as part of the
original pair of assertions, not claimed as a new result.
Coster's classical binomial-power family \cite{Coster} and Straub's
multivariate theorem and weighted cancellation lemma \cite{A:Straub}
are close antecedents; merely having a different displayed sum would not
by itself establish a new theorem or method.

The Baxter walk model and compact-group machinery are classical inputs,
while the specific slice estimate is the analytic contribution offered
for review here. The odd subsequence already follows from the symmetric
walk model by the spectral theorem. The harder even subsequence needs
both domination and mass normalization. The full Baxter sequence is
not claimed to be a Stieltjes moment sequence.

The corpus snapshot used was commit
\texttt{38348353416f1f7fc50cecb93c3b554a00260c13}.
In particular the structural graph proof belongs to VMRC [161], not to
the new transfer argument. Accordingly, neither all-six-gates success
nor superiority to the corpus's earlier results is established by this
manuscript. Its deliverable is the complete, inspectable mathematical
arguments above together with explicit provenance limits.

\begin{thebibliography}{99}
% Verified primary bibliographic records for the Baxter proof.
% Add these items inside the paper's thebibliography environment.
\bibitem{ChungGrahamHoggattKleiman1978}
F. R. K. Chung, R. L. Graham, V. E. Hoggatt, Jr., and M. Kleiman,
\emph{The number of Baxter permutations},
Journal of Combinatorial Theory, Series A \textbf{24} (1978), 382--394.
\href{https://doi.org/10.1016/0097-3165(78)90068-7}{doi:10.1016/0097-3165(78)90068-7}.
Author-hosted primary copy: \url{https://fanchung.ucsd.edu/mypaps/fanpap/24baxter.pdf}.

\bibitem{AsinowskiCardinalFelsnerFusy2025}
A. Asinowski, J. Cardinal, S. Felsner, and \'{E}. Fusy,
\emph{Combinatorics of rectangulations: Old and new bijections},
Combinatorial Theory \textbf{5}(1) (2025).
\href{https://doi.org/10.5070/C65165025}{doi:10.5070/C65165025}.
Preprint: \href{https://arxiv.org/abs/2402.01483}{arXiv:2402.01483} (2024).
Author manuscript: \url{https://igm.univ-mlv.fr/~fusy/Articles/rectangulations.pdf}.
The quadrant-walk representation used here appears on p.~29 of that manuscript.

\bibitem{GuralnikWarnockZemach1985}
G. Guralnik, T. Warnock, and C. Zemach,
\emph{Random selection from trace-biased invariant distributions of
$\mathrm{SU}(3)$ matrices},
Journal of Computational Physics \textbf{61} (1985), 105--118.
\href{https://doi.org/10.1016/0021-9991(85)90063-4}{doi:10.1016/0021-9991(85)90063-4}.
% Its publisher abstract explicitly derives Haar measure in real and imaginary trace
% coordinates. The current proof derives its exact normalization independently;
% no claim that the new slice inequality occurs in this paper is intended.

\bibitem{A:OEIS} OEIS Foundation, A365029, Peter Bala, 27 August 2023.
\url{https://oeis.org/A365029}.
\bibitem{A:Straub} A. Straub, Multivariate Ap\'ery numbers and
supercongruences of rational functions, Algebra Number Theory 8 (2014),
1985--2008; Lemma 5.1. \url{https://arxiv.org/abs/1401.0854}.
\bibitem{A:Beukers} F. Beukers, Some congruences for the Ap\'ery numbers,
J. Number Theory 21 (1985), 141--155.
\bibitem{A:Bajaj} \texttt{rbajaj5}, public research proof portfolio,
A365029FirstTwoLevels.md and A365029AllLevelDigitReduction.md,
consulted 15 November 2026.
\url{https://github.com/rbajaj5/a183068-supercongruence/tree/main/related-results}.

\bibitem{H:DTZ}
H. Deng, Z. Tang and J. Zhang,
On a conjecture of Randi\'c index and graph radius,
\emph{Filomat} \textbf{29} (2015), 1369--1375.
DOI: \href{https://doi.org/10.2298/FIL1506369D}{10.2298/FIL1506369D}.
Preprint \href{https://arxiv.org/abs/1210.2543}{arXiv:1210.2543}, 9 October 2012.
\bibitem{H:Shi}
L. Shi, Chemical indices, mean distance, and radius,
\emph{MATCH Commun. Math. Comput. Chem.} \textbf{75} (2016), 57--70.
\url{https://archive.ymsc.tsinghua.edu.cn/pacm_download/138/4638-match75n1_57-70ChemicalIndex.pdf}.
\bibitem{H:CPS}
M. Cygan, M. Pilipczuk and R. \v{S}krekovski,
On the inequality between radius and Randi\'c index for graphs,
\emph{MATCH Commun. Math. Comput. Chem.} \textbf{67} (2012), 451--466.
\bibitem{H:VMRC161}
Keio and Shaitetsu, A proof of the Caporossi--Hansen conjecture on the
Randi\'c index and graph radius, VMRC manuscript [161], 13 November 2026.
Repository commit \texttt{38348353416f1f7fc50cecb93c3b554a00260c13}.
\url{https://github.com/QEEDDD/Verifiable-Mathematical-Reasoning-Corpus}.

\bibitem{Coster} M. J. Coster, \emph{Supercongruences}, CWI report
AM-R8918 (1989), Theorem 4. \url{https://ir.cwi.nl/pub/5804/5804D.pdf}.
\bibitem{DebCatalogue} B. Deb, \emph{Positivity problems in enumerative and
algebraic combinatorics}, Problems P1.7.5e/o, attributed to S. Dimitrov,
B. Deb and A. D. Sokal. Page last updated 16 June 2026; inspected
15 November 2026.
\url{https://bishaldeb.com/positivity/themes/SMSCHTP/pp-hankel.html}.

\end{thebibliography}
\end{document}
